%------------------------------------------------------------------------------------
%  AVERTISSEMENT PÉDAGOGIQUE ET REMERCIEMENTS
%------------------------------------------------------------------------------------
\addchap{Avertissement Didactique}

\begin{boiteSynthese}{Double Nature de ce Document Modèle}
  Ce volume est conçu selon une double perspective indissociable :
  \begin{enumerate}
    \item \textbf{Côté lecteur :} Il se lit de manière fluide comme un manuel d'ingénierie logicielle et de typographie académique contemporaine.
    \item \textbf{Côté auteur / développeur :} Chaque section est un exemple de code source prêt à l'emploi. Le balisage \LaTeX{} des fichiers sous-jacents a été rédigé avec une clarté exemplaire afin de servir de base de copier-coller immédiate pour vos futurs travaux de thèse, d'articles ou de mémoires de fin d'études.
  \end{enumerate}
\end{boiteSynthese}

\vspace{1cm}

\section*{Remerciements Typographiques et Institutionnels}

Ce projet s'appuie sur la richesse de l'écosystème open-source \LaTeX{} et sur les travaux fondateurs de Donald Knuth et Leslie Lamport, perpétués avec brio par les équipes du projet \LaTeX{}, du moteur Lua\LaTeX{} et de la suite KOMA-Script développée par Markus Kohm.

Nos remerciements s'adressent tout particulièrement :
\begin{itemize}
  \item À la communauté du \emph{Comprehensive \TeX{} Archive Network} (CTAN) pour la maintenance pérenne des packages incontournables mis en œuvre dans cette classe (\textsf{fontspec}, \textsf{unicode-math}, \textsf{tcolorbox}, \textsf{biblatex}, \textsf{cleveref}, \textsf{siunitx}).
  \item Aux enseignants, chercheurs et relecteurs de l'UCLouvain dont les retours méthodologiques ont permis d'éprouver les contraintes de mise en page réelles des mémoires de master et thèses de doctorat.
  \item À l'ensemble des contributeurs aux standards de reproductibilité numérique et d'archivage pérenne des connaissances scientifiques.
\end{itemize}

